\documentclass{article}
\usepackage[T1]{fontenc}
\usepackage{lmodern}
\usepackage{graphicx}
\usepackage{amsmath,amssymb}
\usepackage[a4paper,margin=2cm]{geometry}
\usepackage[absolute,overlay]{textpos}
\setlength{\TPHorizModule}{1mm}
\setlength{\TPVertModule}{1mm}
\usepackage[indonesian]{babel}
\usepackage{hyperref}
\setlength{\emergencystretch}{2em}
% unit: Halaman 28

\begin{document}

% === Halaman 28 (python/native) ===

28

• Bagaimana jika terdapat dua atau lebih nilai k yang menghasilkan pesan-pesan 

bermakna?

  Contoh: Misalkan kriptogram  HSPPW menghasilkan dua kemungkinan  

               kunci yang potensial, yaitu:



 k = 4 menghasilkan pesan dolls     (boneka)



k = 11 menghasilkan wheel  

   (roda)  

. 


  Nilai k mana yang benar?

 Jika kasusnya demikian, maka lakukan dekripsi terhadap potongan cipherteks lain 

tetapi cukup menggunakan k = 4 dan k = 11  agar dapat disimpulkan kunci mana yang 

benar.

\end{document}
